\begin{figure}[tbp]
\centering
\begin{tikzpicture}[x=1cm,y=1cm,font=\small,line cap=round,line join=round]
 \begin{scope}
  \fill[gray!8] (-.35,.3) rectangle (4.35,3.55);
  \draw[thick] (-.35,.3)--(4.35,.3);
  \draw[thick] (-.35,3.55)--(4.35,3.55);
  \node[left] at (-.35,.3) {$S$};
  \node[left] at (-.35,3.55) {$T$};
  \fill[gray!25] (.15,1.12) rectangle (3.9,1.38);
  \fill[gray!25] (.15,2.42) rectangle (3.9,2.68);
  \draw[densely dashed] (-.05,.8) rectangle (4.05,3.0);
  \foreach \x/\i in {.2/0,1.4/1,2.6/2,3.8/3} {
   \draw[thick] (\x,.3)--(\x,3.55);
   \node[above] at (\x,3.55) {$C_{\i}$};
  }
  \draw[-{Stealth}] (.3,1.25)--(3.65,1.25);
  \draw[-{Stealth}] (.3,2.55)--(3.65,2.55);
  \draw (2.6,1.25) circle (.2);
  \draw (2.6,2.55) circle (.2);
  \node[align=center] at (2,-.25) {Two protected openings\\in the common rectangle};
 \end{scope}
 \begin{scope}[xshift=6.2cm]
  \draw[thick] (.1,3.2)--(1.35,2.6)--(2.4,2.6)
       --(3,2.85)--(4.2,2.6);
  \node[above] at (3.5,2.8) {wall};
  \draw[gray] (1.85,1.8) circle (.8);
  \draw[thick,-{Stealth}] (.1,1.25)--(.8,1.25)--(1.85,1.8);
  \draw[thick] (1.85,1.8)--(1.85,2.6);
  \draw[densely dashed] (1.85,1.8)..controls(2.6,1.5)and(3.5,1.65)..(4.2,2.6);
  \fill (1.85,1.8) circle (.035);
  \fill (1.85,2.6) circle (.035);
  \node[below left] at (1.85,1.8) {$x$};
  \node[above] at (1.85,2.6) {$z$};
  \node[right] at (1.85,2.15) {$\delta$};
  \draw[dotted] (1.85,2.6) circle (.2);
  \draw[dotted] (1.85,2.6) circle (.1);
  \node[align=center] at (2,-.25) {Stop at the first distance-$\delta$ contact;\\replace the remaining tail by $xz$};
 \end{scope}
\end{tikzpicture}
\caption{The geometry used to extend an actual crossing.  Left: the
common straight portions of four separated cuts contain two transverse
strips and protected $C_2$ openings.  Outside this rectangle the cuts
may follow irregular walls.  Right: a schematic terminal probe for
$C_3$.  Its retained part is $\delta$-clear, and the nearest-wall segment
from its first distance-$\delta$ point $x$ supplies a tangent approach
to $z$.  The dashed remainder of the old exact-hit tail is discarded;
the dotted circles indicate the smaller terminal neighborhoods.
Neither panel depicts the full disk or asserts separation of distinct
incidences having the same physical boundary position.}
\label{fig:disk-localization}
\end{figure}
